\section{Methods}

In this section, we propose X-codec, a straightforward yet effective method to overcome the semantic shortcomings of the current acoustic codecs. 

\subsection{Acoustic Audio codec}

As illustrated in Figure \ref{fig_xcodec}, our model builds upon the framework established by existing acoustic codecs such as Encodec \cite{defossez2022high} and DAC\cite{kumar2024high}. An acoustic audio codec is composed of three main components: an acoustic encoder, a quantizer, and an acoustic decoder. The input of the codec is the raw waveform $\textbf{X} \in \mathbb{R}^{n}$, where $n$ represents the number of waveform samples. This waveform is fed into the acoustic encoder, which consists of several convolutional layers and employs temporal downscaling to extract frame-level latent acoustic features $\textbf{A} \in \mathbb{R}^{H_a \times T}$, where $H_a$ denotes the hidden size of the acoustic features and $T$ is the number of frames. These continuous features are then transformed into a series of discrete tokens $\textbf{Q} \in \mathbb{R}^{M \times T}$ using a Residual Vector Quantizer (RVQ) with $M$ quantizer layers. During training, a specific codebook for the quantizer is learned, enabling the conversion of discrete tokens back to continuous features $\textbf{A}_q \in \mathbb{R}^{H_a \times T}$. The acoustic decoder then reconstructs the waveform $\hat{\textbf{X}}$ from $\textbf{A}_q$ using several convolutional layers and temporal upsampling. The training process is supervised using various losses, including mel loss, STFT loss, and GAN loss, to ensure high-quality acoustic reconstruction.

\subsection{Analysing Semantic Shortcoming}

In this section, we investigate the impact of acoustic codecs on the performance of audio LLMs, focusing specifically on VALL-E, a pioneering model that leverages language model principles for text-to-speech. Our analysis reveals that training VALL-E using Encodec results in high word error rates (WER) and frequent inaccuracies in content generation. For example, when the input text ``he passed through Henley Saint Albans and came so near to London as Harrow on the Hill'' is synthesized, it is erroneously produced as ``he passed through henley saint albeans and camsel knew to lunglan as herold the lor''. This misinterpretation, which is beyond simply improving the audio quality, suggests a fundamental limitation in Encodec's ability to differentiate phonemes, possibly due to its inadequate semantic processing capabilities.

To substantiate the above hypothesis, we conducted Phonetic Discriminability ABX Tests to evaluate the phonetic discriminability of Encodec's representations. The details are provided in the experiment section. Our findings reveal that Encodec's representations exhibit poor phonetic discriminability, which confirms the presence of semantic inadequacies in the codec. Based on these results, we assert that these semantic shortcomings are a significant contributing factor to the observed inaccuracies of language model based audio generation. 


To effectively address these semantic limitations, we introduce a novel approach that integrates more comprehensive semantic features into the codec's architecture. This enhancement is designed to enrich the codec's understanding of audio content, thereby alleviating the interpreting load on the language model. Detailed elaboration of this method is provided in the subsequent section.


\subsection{Designing Auxiliary Semantic Module}

Our approach employs a straightforward method that enhances audio codecs by directly concatenating semantic and acoustic features. Initially, we extract the semantic feature vector $\textbf{S}^* \in \mathbb{R}^{H_s \times T}$ from the audio waveform $\textbf{x}$. This extraction utilizes a self-supervised, pre-trained model such as HuBERT \cite{hsu2021hubert} or wav2vec 2.0 \cite{baevski2020wav2vec}. The extracted features are then processed through multiple convolutional layers within a semantic encoder to yield the refined semantic feature vector $\textbf{S}$. Concurrently, the acoustic branch produces the feature $\textbf{A}$. These outputs, $\textbf{S}$ and $\textbf{A}$, are subsequently concatenated using a linear projection $\phi$, formulated as:

\begin{equation}
\textbf{U} =  concat(\phi_s(\textbf{S}), \phi_a(\textbf{A})) ,
\end{equation}
where the concatenated feature $\textbf{U} \in \mathbb{R}^{H_u \times T}$ is designed to maximize information preservation from both semantic and acoustic sources. This combined feature is then subject to RVQ using an $M$-layer quantizer, resulting in tokens that encapsulate a rich mixture of semantic and acoustic information.

The quantized feature $\textbf{U}_q$ is designed to meet the decoder's objectives through two projectors, $\beta_s$ and $\beta_a$, which enable the decoders to reconstruct the original semantic feature $\hat{\textbf{S}}^*$ and the audio waveform $\hat{\textbf{x}}$. We adhere to established acoustic reconstruction methods from previous works while introducing a Mean Squared Error (MSE) loss specifically for the reconstruction of semantic features. Furthermore, a constant weight $\gamma$ is applied to the semantic loss to ensure that its scale is aligned with other losses, thus promoting a balanced training objective.